\chapter{Generative AI Use Disclosure}
\label{app:ai-disclosure}

Generative AI tools were used in this project. This appendix discloses each use
in accordance with the aivancity rules: the tool, the period, the task, what was
produced, and what the author verified before inclusion. The source of record is
\texttt{docs/generated\_artifacts\_log.md} in the project repository, an
append-only log of every AI-assisted artefact with its verification status; this
appendix summarises it and does not replace it.

\draftnote{Confirm the tool names, dates and level-of-use wording with the
academic supervisor before submission. The school may require prompts or
transcripts to be attached; the artefact log identifies which files each entry
produced, and the session transcripts can be supplied on request.}

\section{Three distinct kinds of use}

It is important to separate three things that are all, loosely, ``AI use'',
because they carry different obligations.

\paragraph{1. Models under study.} \texttt{gpt-4o-mini-2024-07-18} is the
generator whose behaviour Chapter~\ref{chap:results} measures; the two
\texttt{text-embedding-3} models are components of the retrieval system being
evaluated. These are the object of the research, not tools used to write it.
Every call is recorded in an experiment manifest with its token count, and the
full usage and cost record is in Appendix~\ref{app:reproducibility}. The
historical benchmark questions were themselves generated by a hosted model in the
earlier phase of the project; this is why the benchmark is described throughout
as synthetic in origin and why the independent review in
Section~\ref{sec:review-results} was run.

\paragraph{2. Development assistance.} OpenAI Codex was used in July 2026 for
repository inspection, test and script implementation, experiment-manifest
reconciliation, and LaTeX conversion. Its output was independently re-verified
before acceptance (artefact log, 25 July: ``Independent verification of the
reconciliation''), and that verification found one typesetting regression and
one unconfirmed test count, both recorded.

\paragraph{3. Research and writing assistance.} Claude (Anthropic), accessed
through the Claude desktop application, was used from July to August 2026 as an
assistant for analysis, code, and drafting. This is the most substantial use and
is itemised below.

\section{Itemised record of research and writing assistance}

Table~\ref{tab:ai-use} itemises each AI-assisted artefact, the task it assisted with, and what was verified before it entered this document.

\begin{table}[htbp]
  \centering
  \caption{AI-assisted artefacts, from the artefact log. ``Verified'' states what the author checked before the artefact was accepted.}
  \label{tab:ai-use}
  \footnotesize
  \begin{tabularx}{\textwidth}{@{}l X X@{}}
    \toprule
    \textbf{Date} & \textbf{Produced with AI assistance} & \textbf{Verified by the author} \\
    \midrule
    25 Jul & Review-form generator, response aggregator, pseudonymisation script & Round-tripped against synthetic responses; a defect in the $\kappa$ computation was found and fixed \\
    25 Jul & Chapter 2 draft; bibliography from 10 to 64 entries & Every entry checked against a primary record; the chapter marked pending author edit \\
    25 Jul & Chapter 4 draft & Every quantity traced to a manifest; eight repository discrepancies recorded rather than smoothed \\
    27 Jul & Chapter 1 and 3 drafts & Regulatory content read from the validated corpus, no provision asserted from memory \\
    8 Aug & Adjudication tool; freeze; retrieval re-scoring; corrected citation evaluator with 29 tests & Contract validated; a split-recomputation defect caught and reversed (D028); the old scorer's errors pinned by tests \\
    8 Aug & Generation runner; verification rule; prompt v2; escalation calibration & Held-out split protected by tooling; every result compared with its artefact; one prose claim corrected by the artefact (see log) \\
    8--15 Aug & Chapters 5 and 6, abstract, appendices; consistency corrections to Chapters 1--4 & Every figure read from a stored metrics file; the two properties no internal metric detects checked independently \\
    \bottomrule
  \end{tabularx}
\end{table}

\section{What was not delegated}

The following were decided by the author and not by any AI system.

\begin{itemize}
  \item The research question and its reframing from a system comparison into a
    failure analysis.
  \item Every one of the 44 recorded decisions in the decision log, including
    the choice to run an independent review, the pre-registered cutoff, the
    choice of generator, and the decision not to fit an escalation threshold.
  \item The adjudication of the 19 contested benchmark records, each with a
    recorded justification.
  \item The authorisation of every hosted API call.
  \item Whether each AI-proposed correction was accepted. Three proposed
    adjudication repairs (D026) were approved by the author before application,
    with the pre-repair file retained.
\end{itemize}

\section{Provenance discipline}

Two rules governed AI-assisted work throughout. Every number in this document
comes from a stored, hashed artefact rather than from prose, so an AI-drafted
sentence cannot introduce a figure that does not exist in the record. And where
the artefact and the prose disagreed, the artefact was allowed to overrule the
prose; the artefact log records the occasions on which that happened.

The chapters were drafted with assistance and are the author's responsibility in
full. The author has read, and where necessary rewritten, every chapter, and is
prepared to defend every claim in them orally.
